% ============================================================
%  EXAMEN TEMA 3: TRANSFORMACIÓN DE UNIDADES Y VECTORES
%  Curso 4to Año — Física
%  Duración: 90 minutos
%  Temas: Sistema Internacional, transformación de unidades y
%         vectores (componentes, suma y reducción al 1er cuadrante)
%  Autor: Ing. Gabriel Astudillo
%  Nota: enunciado limpio (sin pistas ni desarrollo). Las
%  soluciones están en la "Clave de Respuestas" al final.
% ============================================================

\documentclass[12pt, letterpaper]{article}

% ── Paquetes ─────────────────────────────────────────────────

\usepackage{mathwizards-guia}
\mwguia{Examen — Tema 3: Unidades y Vectores}{Ing. Gabriel Astudillo $\cdot$ 4to Año}

% ── Comandos propios de este examen ───────────────────
\providecommand{\sen}{\sin}

\begin{document}
% ============================================================

% ── Portada / Título ─────────────────────────────────────────
\mwportada{Examen del Tema 3}{Transformación de Unidades e Introducción a los Vectores}{Sistema Internacional, Factores de Conversión y Magnitudes con Dirección}

\vspace{4pt}
\begin{cuadroinfo}
  \small
  \textbf{Nombre:} \rule{0.55\linewidth}{0.4pt} \quad \textbf{Fecha:} \rule{0.15\linewidth}{0.4pt} \\[6pt]
  \textbf{Tiempo disponible:} 90 minutos \quad$\bullet$\quad \textbf{Puntaje total:} 100 puntos \\[4pt]
  \textbf{Instrucciones:}
  \begin{enumerate}[leftmargin=*]
    \item Lee cada ejercicio con calma antes de resolverlo.
    \item Muestra \emph{todos} tus pasos, incluidos los factores de conversión.
    \item \textbf{No olvides las unidades} en tus resultados.
    \item Usa los valores exactos de los ángulos notables cuando aparezcan.
    \item No se permite el uso de calculadora.
    \item Escribe con letra clara y revisa tu examen antes de entregarlo.
  \end{enumerate}
\end{cuadroinfo}

\vspace{8pt}

% ============================================================
\section{Sistema Internacional de Unidades y Medición (28 puntos)}
% ============================================================

\begin{enumerate}[leftmargin=*]
  \item \textbf{(8 pts)} Clasifica cada magnitud como fundamental o derivada e indica su unidad del SI:
        \[
          a)\ \text{longitud} \qquad b)\ \text{velocidad} \qquad c)\ \text{temperatura} \qquad d)\ \text{fuerza} \qquad e)\ \text{tiempo}
        \]

  \item \textbf{(8 pts)} Expresa en notación científica:
        \[
          a)\ 0{,}000\,000\,36 \qquad b)\ 78\,000\,000 \qquad c)\ 0{,}000\,42
        \]

  \item \textbf{(6 pts)} Indica cuántas cifras significativas tiene cada medida:
        \[
          a)\ 2{,}50\;\text{m} \qquad b)\ 0{,}004\;\text{kg} \qquad c)\ 3{,}100\;\text{s}
        \]

  \item \textbf{(6 pts)} \textquestiondown Cuántos milímetros hay en $3{,}2\;\text{m}$? \textquestiondown Y cuántos centímetros?
\end{enumerate}

\newpage

% ============================================================
\section{Transformación de Unidades (34 puntos)}
% ============================================================

\begin{enumerate}[leftmargin=*]
  \item \textbf{(10 pts)} Convierte usando factores de conversión:
        \[
          a)\ 350\;\text{cm} \to \text{m} \qquad b)\ 4{,}5\;\text{km} \to \text{m} \qquad c)\ 2\,500\;\text{g} \to \text{kg} \qquad d)\ 2\;\text{h} \to \text{s} \qquad e)\ 540\;\text{km/h} \to \text{m/s}
        \]

  \item \textbf{(8 pts)} Convierte:
        \[
          a)\ 15\;\text{m/s} \to \text{km/h} \qquad b)\ 108\;\text{km/h} \to \text{m/s}
        \]

  \item \textbf{(8 pts)} Un bloque tiene una masa de $300\;\text{g}$ y ocupa un volumen de $40\;\text{cm}^3$. Calcula su densidad en $\text{g/cm}^3$ y exprésala también en $\text{kg/m}^3$.

  \item \textbf{(8 pts)}
        \begin{enumerate}[label=\alph*)]
          \item Una maratón mide $42{,}195\;\text{km}$. \textquestiondown Cuántos metros son?
          \item Un tren viaja a $90\;\text{km/h}$. \textquestiondown Cuántos metros recorre en un segundo?
        \end{enumerate}
\end{enumerate}

\newpage

% ============================================================
\section{Vectores (38 puntos)}
% ============================================================

\begin{enumerate}[leftmargin=*]
  \item \textbf{(8 pts)} Clasifica cada magnitud como escalar (E) o vectorial (V):
        \[
          a)\ \text{masa} \qquad b)\ \text{desplazamiento} \qquad c)\ \text{velocidad} \qquad d)\ \text{temperatura}
        \]

  \item \textbf{(10 pts)} Un vector $\vec{v}$ tiene módulo $10$ y forma un ángulo de $30^{\circ}$ con el eje $X$ positivo. Halla sus componentes rectangulares usando los valores exactos.

  \item \textbf{(10 pts)} Suma los vectores $\vec{a} = (6, -2)$ y $\vec{b} = (-2, 5)$.
        \begin{enumerate}[label=\alph*)]
          \item Escribe el vector resultante como par ordenado.
          \item Calcula el módulo de la resultante.
        \end{enumerate}

  \item \textbf{(10 pts)} Un vector $\vec{u}$ tiene módulo $12$ y forma un ángulo de $210^{\circ}$ con el eje $X$ positivo.
        \begin{enumerate}[label=\alph*)]
          \item Indica el cuadrante y el ángulo de referencia.
          \item Halla sus componentes $u_x$ y $u_y$.
          \item Verifica que $\sqrt{u_x^2 + u_y^2} = 12$.
        \end{enumerate}
\end{enumerate}

\newpage

% ============================================================
\ifmwclaves
  \section{Clave de Respuestas}
  % ============================================================

  \subsection*{Sección 1: Sistema Internacional de Unidades y Medición}

  \begin{enumerate}[leftmargin=*, label=\textbf{\arabic*.}]
    \item a) Fundamental, metro (m). \quad b) Derivada, m/s. \quad c) Fundamental, kelvin (K). \quad d) Derivada, newton (N). \quad e) Fundamental, segundo (s).

    \item a) $3{,}6 \times 10^{-7}$ \quad b) $7{,}8 \times 10^{7}$ \quad c) $4{,}2 \times 10^{-4}$.

    \item a) 3 cifras. \quad b) 1 cifra (los ceros a la izquierda no cuentan). \quad c) 4 cifras (los ceros a la derecha sí cuentan).

    \item $3{,}2\;\text{m} = 3\,200\;\text{mm}$ y $3{,}2\;\text{m} = 320\;\text{cm}$.
  \end{enumerate}

  \newpage

  \subsection*{Sección 2: Transformación de Unidades}

  \begin{enumerate}[leftmargin=*, label=\textbf{\arabic*.}]
    \item a) $350 \div 100 = 3{,}5\;\text{m}$.
          \quad b) $4{,}5 \times 1\,000 = 4\,500\;\text{m}$.
          \quad c) $2\,500 \div 1\,000 = 2{,}5\;\text{kg}$.
          \quad d) $2 \times 3\,600 = 7\,200\;\text{s}$.
          \quad e) $540 \div 3{,}6 = 150\;\text{m/s}$.

    \item a) $15 \times 3{,}6 = 54\;\text{km/h}$. \quad b) $108 \div 3{,}6 = 30\;\text{m/s}$.

    \item $\rho = \dfrac{300}{40} = 7{,}5\;\text{g/cm}^3$. \quad En $\text{kg/m}^3$: $7{,}5 \times 1\,000 = 7\,500\;\text{kg/m}^3$.

    \item a) $42{,}195\;\text{km} = 42\,195\;\text{m}$.
          \quad b) $90 \div 3{,}6 = 25\;\text{m/s}$, así que recorre $25\;\text{m}$ en un segundo.
  \end{enumerate}

  \newpage

  \subsection*{Sección 3: Vectores}

  \begin{enumerate}[leftmargin=*, label=\textbf{\arabic*.}]
    \item a) E \quad b) V \quad c) V \quad d) E

    \item Con $\cos 30^{\circ} = \dfrac{\sqrt{3}}{2}$ y $\sen 30^{\circ} = \dfrac{1}{2}$:
          \[
            v_x = 10\cos 30^{\circ} = 10 \cdot \frac{\sqrt{3}}{2} = 5\sqrt{3} \approx 8{,}66, \qquad
            v_y = 10\sen 30^{\circ} = 10 \cdot \frac{1}{2} = 5.
          \]
          Entonces $\vec{v} = (5\sqrt{3},\, 5)$.

    \item a) $\vec{a} + \vec{b} = (6 + (-2),\, -2 + 5) = (4, 3)$.
          \quad b) $|\vec{a} + \vec{b}| = \sqrt{4^2 + 3^2} = \sqrt{16 + 9} = \sqrt{25} = 5$.

    \item a) $210^{\circ}$ está en el \textbf{cuadrante III} y el ángulo de referencia es $\theta_{\text{ref}} = 210^{\circ} - 180^{\circ} = 30^{\circ}$.
          \quad b) Ambas componentes son negativas:
          \[
            u_x = -12\cos 30^{\circ} = -12 \cdot \frac{\sqrt{3}}{2} = -6\sqrt{3} \approx -10{,}39, \qquad
            u_y = -12\sen 30^{\circ} = -12 \cdot \frac{1}{2} = -6.
          \]
          \quad c) $\sqrt{(-6\sqrt{3})^2 + (-6)^2} = \sqrt{108 + 36} = \sqrt{144} = 12$. \ \checkmark
  \end{enumerate}

\fi
\end{document}
